\section{The natural quotient frame}

The canonical bipartite double cover carries an evident regular two-fold
covering structure. In this section we develop this distinguished quotient
frame and determine its stabilizer inside the automorphism group.

\subsection{Natural fibers and the deck involution}

Recall that

\[
X=\KC(L(P))
\]

has vertex set

\[
V(X)
=
V(L(P))\times\Ftwo.
\]

For each

\[
u\in V(L(P))
\]

define the fiber

\[
F_u
=
\{
(u,0),
(u,1)
\}.
\]

The collection

\[
\mathcal P_0
=
\{
F_u
:
u\in V(L(P))
\}
\]

is a partition of

\[
V(X)
\]

into fifteen unordered pairs.

The associated involution is

\[
\tau_0(u,\varepsilon)
=
(u,\varepsilon+1),
\]

which exchanges the two vertices of every fiber.

Since

\[
\tau_0^2=\mathrm{id},
\]

and

\[
\tau_0
\]

preserves adjacency, it is a fixed-point-free involutory automorphism of
\(X\).

The partition

\[
\mathcal P_0
\]

is therefore a regular two-fold quotient frame in the sense of
Definition~2.3.


\subsection{The natural quotient}

The quotient graph determined by the fiber partition

\[
\mathcal P_0
\]

is naturally identified with the line graph of the Petersen graph.

\begin{proposition}

The quotient projection

\[
\pi:
X
\longrightarrow
X/\mathcal P_0
\]

is a regular two-fold covering projection and

\[
X/\mathcal P_0
\cong
L(P).
\]

\end{proposition}

\begin{proof}

Every block of

\[
\mathcal P_0
\]

is a fiber

\[
\{
(u,0),
(u,1)
\}
\]

over a unique vertex

\[
u\in L(P).
\]

Define

\[
\Phi:
X/\mathcal P_0
\longrightarrow
L(P)
\]

by

\[
\Phi(F_u)=u.
\]

This map is clearly bijective.

Two quotient vertices

\[
F_u,
\qquad
F_v
\]

are adjacent precisely when some vertex of

\[
F_u
\]

is adjacent to some vertex of

\[
F_v.
\]

By construction of the canonical bipartite double cover,

\[
(u,\varepsilon)
\sim
(v,\varepsilon+1)
\]

if and only if

\[
u\sim v
\]

in

\[
L(P).
\]

Hence

\[
F_u
\sim
F_v
\]

in the quotient exactly when

\[
u
\sim
v
\]

in

\[
L(P).
\]

Therefore

\[
\Phi
\]

is a graph isomorphism.

Since every edge of

\[
L(P)
\]

lifts to two disjoint edges joining the corresponding fibers, the quotient
projection is a regular two-fold covering projection.

\end{proof}

Consequently, the fiber partition

\[
\mathcal P_0
\]

recovers the original base graph from which the canonical bipartite double
cover was constructed.


\subsection{The lifted \(S_{5}\)-action}

Every automorphism

\[
g\in\Aut(L(P))
\]

lifts naturally to an automorphism of the canonical bipartite double cover by

\[
\widehat g(u,\varepsilon)
=
(g(u),\varepsilon).
\]

Since adjacency in

\[
X=\KC(L(P))
\]

depends only on adjacency in the base graph together with the parity
coordinate, every lifted map preserves adjacency.

Moreover,

\[
\widehat g
\]

commutes with the deck involution

\[
\tau_0,
\]

since

\[
\widehat g\,\tau_0
=
\tau_0\,\widehat g.
\]

Thus the lifted action preserves every fiber of the natural quotient frame
\(\mathcal P_0\).

\begin{theorem}

The stabilizer of the natural quotient frame is

\[
A_{\mathcal P_0}
\cong
S_{5}\times C_{2}.
\]

\end{theorem}

\begin{proof}

Every automorphism of the base graph has a unique sheet-preserving lift
to the cover. These lifts give an embedded copy of

\[
S_{5}
\cong
\Aut(L(P)).
\]

Each base automorphism also has a second frame-preserving lift, obtained
by composing its sheet-preserving lift with the deck involution

\[
\tau_0.
\]

The deck involution itself preserves every fiber.

Since every lifted automorphism commutes with

\[
\tau_0,
\]

these actions generate a subgroup

\[
S_{5}\times C_{2}.
\]

Conversely, let

\[
\varphi
\]

preserve the quotient frame

\[
\mathcal P_0.
\]

Then

\[
\varphi
\]

induces an automorphism of the quotient graph

\[
X/\mathcal P_0
\cong
L(P),
\]

whose lift differs from

\[
\varphi
\]

by at most the deck involution.

Hence every frame-preserving automorphism lies in

\[
S_{5}\times\langle\tau_0\rangle,
\]

proving

\[
A_{\mathcal P_0}
\cong
S_{5}\times C_{2}.
\]

\end{proof}

Theorem~4.2 describes the symmetry group visible from the natural quotient
frame. Section~6 shows that the full automorphism group is strictly larger,
arising from additional symmetries that exchange quotient frames rather
than preserve a single one.

